% Series cover adapted from AI Agent Book. Pure TikZ: compute, fabric, cluster.
\begin{titlepage}
\thispagestyle{empty}
\begin{tikzpicture}[remember picture,overlay]
  \fill[structurecolor] (current page.north west) rectangle ([yshift=-0.85cm]current page.north east);
  \fill[structurecolor] (current page.south west) rectangle ([yshift=0.5cm]current page.south east);
\end{tikzpicture}
\centering
\vspace*{2.5cm}
{\fontsize{37}{46}\selectfont\rmfamily\bfseries 深入理解 AI Infra\par}
\vspace{0.5cm}
{\Large\sffamily\color{structurecolor} 量化分析与系统设计\par}
\vspace{1.6cm}
\begingroup
\begin{tikzpicture}[line join=round,line cap=round,
  ink/.style={draw=structurecolor,line width=1.05pt},
  module/.style={ink,fill=white,rounded corners=2pt},
  wire/.style={ink,line width=0.9pt}]
  \path[use as bounding box] (-3.7,-3.5) rectangle (3.7,3.6);
  % One processor: a grid of compute tiles, a shared rim and external pins.
  \draw[module,fill=structurecolor!5] (-.96,1.12) rectangle (.96,3.04);
  \draw[ink,rounded corners=1pt] (-.65,1.43) rectangle (.65,2.73);
  \foreach \x in {-.4,0,.4}{
    \foreach \y in {1.68,2.08,2.48}{
      \fill[structurecolor,rounded corners=.6pt] (\x-.105,\y-.105) rectangle (\x+.105,\y+.105);
    }
    \draw[wire] (\x,3.04)--(\x,3.28);
    \draw[wire] (\x,1.12)--(\x,.88);
  }
  \foreach \y in {1.68,2.08,2.48}{
    \draw[wire] (-1.2,\y)--(-.96,\y);
    \draw[wire] (.96,\y)--(1.2,\y);
  }
  % Fabric: two parallel paths join the processor to the cluster.
  \draw[wire] (-.4,.88)--(-.4,.30);
  \draw[wire] (.4,.88)--(.4,.30);
  \draw[module,fill=structurecolor!5] (-1.38,-.30) rectangle (1.38,.30);
  \foreach \x in {-.95,-.57,-.19,.19,.57,.95}{
    \draw[ink] (\x-.09,-.09) rectangle (\x+.09,.09);
  }
  \draw[wire] (0,-.30)--(0,-.82);
  \draw[wire] (-2.45,-1.35)--(-2.45,-.82)--(2.45,-.82)--(2.45,-1.35);
  \draw[wire] (0,-.82)--(0,-1.35);
  \fill[structurecolor] (0,-.82) circle (.045);
  % Three racks, each with three nodes. Repetition conveys scale.
  \foreach \x in {-2.45,0,2.45}{
    \draw[module] (\x-.79,-3.25) rectangle (\x+.79,-1.35);
    \foreach \y in {-1.73,-2.27,-2.81}{
      \draw[module,fill=structurecolor!4] (\x-.57,\y-.19) rectangle (\x+.57,\y+.19);
      \fill[structurecolor] (\x-.37,\y) circle (.035);
      \draw[wire] (\x-.15,\y)--(\x+.34,\y);
    }
  }
\end{tikzpicture}
\endgroup
\vfill
{\LARGE\sffamily 李博杰\quad 著\par}
\vspace{0.7cm}
{\small\sffamily\color{structurecolor!80} \BookEdition\ · \number\year 年 \number\month 月 \number\day 日\par}
\vspace{1.2cm}
\end{titlepage}
